\section{Gr\"obner bases as confluent reduction systems}\label{sec:groebner-criterion}

Definition~\ref{def:groebner-conditions} distinguishes the self-relative
predicate
\leandoc{Mathlib/RingTheory/MvPolynomial/GroebnerBasisCriterion}{MonomialOrder.IsGroebner}{MonomialOrder.IsGroebner}
from the ideal-relative predicate
\leandoc{Mathlib/RingTheory/MvPolynomial/BuchbergerCriterion}{MonomialOrder.IsGroebnerBasis}{MonomialOrder.IsGroebnerBasis}.
The theorem
\leandoc{Mathlib/RingTheory/MvPolynomial/BuchbergerCriterion}{MonomialOrder.isGroebnerBasis_span_iff_isGroebner}{isGroebnerBasis_span_iff_isGroebner}
identifies them when
$I=\ideal{G}$.  Listing~\ref{lst:groebner-defs} displays the self-relative
Lean definition and its confluence characterization.

\begin{lstlisting}[style=lean,caption={The self-relative Gr\"obner condition and its confluence characterization.},label={lst:groebner-defs}]
def IsGroebner
    (G : Set (MvPolynomial σ R)) : Prop :=
  Ideal.span
      (m.leadingTerm ''
        (((Ideal.span G : Ideal (MvPolynomial σ R)) :
          Set (MvPolynomial σ R)))) =
    Ideal.span (m.leadingTerm '' G)

theorem IsGroebner.iff_confluent
    (G : Set (MvPolynomial σ R))
    (hG₀ :
      ∀ g ∈ G,
        IsUnit (m.leadingCoeff g) ∨ g = 0) :
    m.IsGroebner G ↔
      Relation.Confluent (m.ReducesToSet G)
\end{lstlisting}

\begin{theorem}[Reduction-theoretic characterizations of Gr\"obner bases]
\label{thm:groebner-characterizations}
Assume Condition~\eqref{eq:unit-condition}.  Then the following propositions
are equivalent:
\begin{enumerate}[label=(\alph*),leftmargin=2em]
  \item $\ideal{\LT(f)\mid f\in\ideal{G}}
        =\ideal{\LT(g)\mid g\in G}$;
  \item every $f\in\ideal{G}$ satisfies $f\redstar{G}0$;
  \item every nonzero element of $\ideal{G}$ is reducible modulo $G$;
  \item $\red{G}$ satisfies the Church--Rosser property;
  \item $\red{G}$ is locally confluent;
  \item $\red{G}$ is confluent;
  \item $\red{G}$ has unique normal forms; and
  \item each residue class modulo $\ideal{G}$ has a unique normal-form
        representative.
\end{enumerate}
\end{theorem}

\begin{table}[t]
\caption{Theorem~\ref{thm:groebner-characterizations} and its Lean declarations.}
\label{tab:groebner-characterizations-lean}
\centering
\scriptsize
\begin{tabularx}{\textwidth}{@{}>{\raggedright\arraybackslash}p{0.07\textwidth}>{\raggedright\arraybackslash}p{0.43\textwidth}>{\raggedright\arraybackslash}X@{}}
\toprule
Label & Mathematical condition & Lean declaration \doclinkmark \\
\midrule
(a) & Leading-term ideal criterion & \leandoc{Mathlib/RingTheory/MvPolynomial/GroebnerBasisCriterion}{MonomialOrder.IsGroebner}{MonomialOrder.IsGroebner} \\
(b) & Every element of $\ideal{G}$ reduces to zero & \leandoc{Mathlib/RingTheory/MvPolynomial/GroebnerBasisCriterion}{MonomialOrder.IsGroebner.iff_idealElements_reduceTo_zero}{iff_idealElements_reduceTo_zero} \\
(c) & Every nonzero element of $\ideal{G}$ is reducible & \leandoc{Mathlib/RingTheory/MvPolynomial/GroebnerBasisCriterion}{MonomialOrder.IsGroebner.iff_nonzero_idealElements_reducible}{iff_nonzero_idealElements_reducible} \\
(d) & Church--Rosser & \leandoc{Mathlib/RingTheory/MvPolynomial/GroebnerBasisCriterion}{MonomialOrder.IsGroebner.iff_churchRosser}{iff_churchRosser} \\
(e) & Local confluence & \leandoc{Mathlib/RingTheory/MvPolynomial/GroebnerBasisCriterion}{MonomialOrder.IsGroebner.iff_locallyConfluent}{iff_locallyConfluent} \\
(f) & Confluence & \leandoc{Mathlib/RingTheory/MvPolynomial/GroebnerBasisCriterion}{MonomialOrder.IsGroebner.iff_confluent}{iff_confluent} \\
(g) & Unique normal forms & \leandoc{Mathlib/RingTheory/MvPolynomial/GroebnerBasisCriterion}{MonomialOrder.IsGroebner.iff_uniqueNormalForms}{iff_uniqueNormalForms} \\
(h) & Unique normal-form representative in each residue class & \leandoc{Mathlib/RingTheory/MvPolynomial/GroebnerBasisCriterion}{MonomialOrder.IsGroebner.iff_unique_normalForm_representatives}{iff_unique_normalForm_representatives} \\
\bottomrule
\end{tabularx}
\medskip

\raggedright\scriptsize
The declarations corresponding to (b)--(h) are in the namespace
\lean{MonomialOrder.IsGroebner}.
\end{table}

The proof of Theorem~\ref{thm:groebner-characterizations} is organized as
explicit implications between the eight propositions.  Termination gives at
least one normal form reachable from every polynomial.  The equality in
(a) proves that a residue class modulo $\ideal{G}$ contains at most one
normal polynomial, and hence exactly one.  By Theorem~\ref{thm:translation},
equality of residue classes is the same as reduction equivalence, so uniqueness
of normal representatives yields Church--Rosser.
Theorem~\ref{thm:reduction-termination} and Newman's lemma identify local
confluence with confluence.  The Church--Rosser property implies that every
$f\in\ideal{G}$ reduces to zero because $f\eqvred{G}0$.

The converse from (c) to (a) requires more than an immediate
leading-term argument: a nonzero ideal element may initially be reducible at a
lower term.  The Lean proof uses well-founded induction on the support-colex
order.  If the reduction occurs below the leading term, the reduct has the same
leading term and a strictly smaller support-colex measure, so the induction
hypothesis applies.  Eventually one obtains a reduction at the leading term,
which shows that the leading term is divisible by the leading term of an
element of $G$ and proves the equality in (a).
